\subsection{Moves for Inequalities} 
We have the following moves for solving equations:
\begin{enumerate}
\item Basic moves (add, subtract, multiply, divide)
\item Treat as a polynomial and factor
\item Rewrite using an inverse
\end{enumerate}
We've seen how the first two apply to inequalities, with some warnings.
\begin{enumerate}
\item Basic moves---multiplying by a \makeblank[1in] number reverses the inequality.
\item Factoring polynomials---detailed sign analysis needed
\item Rewrite using an inverse---?
\end{enumerate}
When we solved \makeblank inequalities, we had to think about the \makeblank, because the \makeblank function can be \emph{increasing} or \emph{decreasing}.
\begin{definition}
A function $f$ is \term{increasing} if $f(a)>f(b)$ if and only if $a>b$. A function is \term{decreasing} if $f(a)<f(b)$ whenever $a>b$.
\end{definition}
\begin{Example}
$f(x)=3x$ is an increasing function, because if $a>b$, then $3a>3b$ and vice versa.
\end{Example}
\begin{Example}
$g(x)=-2x$ is a decreasing function, because if $a>b$, then $-2a<-2b$ and vice versa.
\end{Example}
This gives us the third move for solving inequalities.
\begin{itemize}
\item We can apply any function to both sides of an inequality.
\item If the function is \makeblank[1in], the inequality sign \makeblank[1.5in]
\item If the function is \makeblank[1in], the inequality sign \makeblank[1.5in]
\item If the function does both, we must \makeblank as we did the absolute value.
\end{itemize}
The $\log_b$ is an increasing function if $b\makeblanksmall$, and decreasing if $b\makeblanksmall$.